% T12 C/D, chapitre 12 : Primitives et intégrales. Source : le livre « Terminales scientifiques » (Integrale/Lesona, Integrale/Exo).
% Les exercices tirés d’autres ouvrages ne sont pas repris (décision du 2 novembre 2026).

\TCDLessonStyleOn
\providecommand{\dx}{~\mathrm{d}x}
\providecommand{\dt}{~\mathrm{d}t}

\chapter{Primitives et Int\'egrales}
      \section{Primitives}
      \begin{defi}[Primitive]
      \index{Primitives}
            Soit $f$ une fonction d\'efinie sur un intervalle $I$ non-vide.
            Une fonction $F$ est une \emph{primitive} de $f$ sur $I$ si $F$ est d\'erivable sur $I$ et si, 
            pour tout $x\in I$, $F'(x)=f(x)$.

      \smallnota{%
            Si une fonction $f$ admet une primitive $F$, alors elle en admet une infinit\'e.
            En fait, toutes les fonctions de la forme $F+c$ o\`u $c\in\mathbb{R}$ sont aussi des primitives de $f$.
            }
            
            L'ensemble des primitives de $f$ est not\'e par $\displaystyle{\int f(x) \dx}$.
            \begin{ex}~\\
                  \begin{itemize}
                        \item $f:x\mapsto 2x$ admet $F:x\mapsto x^2$ comme primitive sur $\mathbb{R}$.
                        \item $g:x\mapsto \frac{1}{2\sqrt{x}}$ admet $G:x\mapsto \sqrt{x} + 2$ 
                              comme primitive sur $]0,+\infty[$.
                  \end{itemize}
            \end{ex}
      \end{defi}
      %
      \begin{thm}
      Soit $f$ une fonction d\'efinie sur un intervalle $I$ et $F$ une primitive de $f$ sur $I$.
      \begin{enumerate}
            \item 
                  La fonction $\lambda f$ admet la fonction $\lambda F$ comme primitive sur $I$.
            \item 
                  Soit $g$ une fonction d\'efinie sur un intervalle $I$ et $G$ une primitive de $g$ sur $I$.
                  Alors la fonction $f+g$ admet la fonction $F+G$ comme primitive sur $I$.      
            \item 
                  Soit $F_1$ une autre primitive de $f$ sur $I$.
                  Alors la fonction $F - F_1$ est constante.
            \item 
                  Soient $x_0\in I$ et $y_0\in \mathbb{R}$.
                  Alors $f$ admet une unique primitive $F_0$ sur $I$ v\'erifiant $F_0 (x_0) = y_0$.
      \end{enumerate}
            \begin{ex}
                  Cherchons la primitive de la fonction $f:x\mapsto 3 x^2$ qui s'annule en $1$.
                  La fonction $F:x\mapsto x^3$ est d\'erivable et sa d\'eriv\'ee est $F':x\mapsto 3x^2$.
                  Comme $F(1)=1$, alors on peut prendre $F_0 : x\mapsto x^3 - 1$.
                  Ainsi, $F_0' (x) = 3x^2$ et $F_0' (1) = 1 - 1 = 0$.
            \end{ex}
      \end{thm}
%
      \begin{prop}[Primitives usuelles] ~\\
            \begin{minipage}{0.45\actualtextwidth}
                  \begin{tabular}{|c|c|}
                        \hline                                                            %\raisebox{-6pt}{$\rule{0pt}{17pt}$}
                        \headingstyleko
                        Fonctions                   & Primitives                \\ \hline %\raisebox{-6pt}{$\rule{0pt}{17pt}$}
                        $0$                         & $k$                       \\ \hline %\raisebox{-6pt}{$\rule{0pt}{17pt}$}
                        $1$                         & $x + k$                   \\ \hline %\raisebox{-6pt}{$\rule{0pt}{17pt}$}
                        $\frac{1}{x^2}$             & $-\frac{1}{x}+k$          \\ \hline %\raisebox{-6pt}{$\rule{0pt}{17pt}$}
                        $\frac{1}{\sqrt{x}}$        & $2\sqrt{x}+k$             \\ \hline %\raisebox{-6pt}{$\rule{0pt}{17pt}$}
                        $x^n \quad (n\neq -1)$      & $\frac{x^{n+1}}{n+1} + k$ \\ \hline %\raisebox{-6pt}{$\rule{0pt}{17pt}$}
                        $\cos x$                    & $\sin x + k$              \\ \hline %\raisebox{-6pt}{$\rule{0pt}{17pt}$}
                        $\sin x$                    & $-\cos x +k$              \\ \hline %\raisebox{-6pt}{$\rule{0pt}{17pt}$}
                        $\expo{x}$                  & $\expo{x} + k$            \\ \hline %\raisebox{-6pt}{$\rule{0pt}{17pt}$}
                        $\frac{1}{x}$               & $\ln |x| + k$             \\ \hline 
                  \end{tabular}            
            \end{minipage}
      %
            \begin{minipage}{0.45\actualtextwidth}
                  Notons que la recherche de primitives est en g\'en\'erale difficile.
                  Bien que toute fonction continue admette une primitive, 
                  il existe des fonctions dont on ne 
                  sait pas donner des primitives sous forme de fonctions usuelles.

                  En voici quelques exemples :
						\begin{itemize}
							\item $x \mapsto \expo{-x^2}$
							\item $x \mapsto \sqrt{x^4 + 1}$
							\item $x \mapsto (\ln x)^{-1}$
							\item $x \mapsto \sqrt{x} \cdot \sin (x) $
							\item \ldots
						\end{itemize}
            \end{minipage}
      \end{prop}
      %
      \begin{thm}
            Soient $f$ et $g$ deux fonctions d\'erivables sur un intervalle $I$.
            Comme la primitive est la r\'eciproque de la d\'erivation, on a :
                  \begin{IEEEeqnarray*}{rCl}
                        \int g'(x) \cdot f'\left( g(x) \right) \dx & = & f\left( g(x) \right) + C       \\
                        \int f'(x) \cdot f^n (x) \dx               & = & \frac{1}{n+1} f^{n+1}(x) + C
                  \end{IEEEeqnarray*}
            %
            \begin{ex}~\vspace{-1\baselineskip}
                  \begin{enumerate}
                        \item On veut trouver une primitive de $a(x) = 2x \cdot \cos (x^2)$.
                              On constate que $2x$ est en fait la d\'eriv\'ee de $x^2$.
                              Cette fonction est donc de la forme $u'(x) \cdot \cos u(x)$.
                              Il nous reste donc \`a trouver la primitive de la fonction $\cos$,
                              qui n'est autre que la fonction $\sin$.
                              Ainsi, 
                                    \begin{align*}
                                    a(x) & = (x^2)' \cdot \sin' (x^2)            \\
                                         & = g'(x) \cdot f'\left(g(x)\right)     \,.
                                    \end{align*}
                              Ainsi, la fonction $A$ d\'efinie par $A(x) = \sin x^2$ est une primitive de $a$.
                        \item On veut trouver une primitive de $b(x) = \cos x \cdot \sin^4 x$.
                              On commence par remarquer que $\cos$ est en fait la d\'eriv\'ee de $\sin$.
                              Ainsi, $b(x) = \sin' x \cdot \sin^4 x$.
                              La fonction $B$ d\'efinie par $B(x) = \frac{\sin^5 x}{5}$ est une primitive de $b$.
                  \end{enumerate}
            \end{ex}
            %
      \end{thm}
%%%%%%%%%%%%%%%%%%%%%%%%%%%%
      \section{Int\'egrales}
      \subsection{G\'en\'eralit\'es}
	      \begin{defi}
	      \index{Integrales@Int\'egrales}
	      \index{Integration@Int\'egration}
	      Soit $f$ une fonction d\'efinie sur un intervalle $I$ et y admettant une primitive $F$.
	      Soient $a$ et $b$ deux \'el\'ements de $I$.
	      On appelle \emph{int\'egrale de $f$} entre $a$ et $b$ le r\'eel 
	            \[
	            \int_{a}^{b} f(x) \dx = \left[ F(x) \right]_{a}^{b} = F(b) - F(a)      \,.
	            \]
	      \warning{%
	              La variable $x$ de la d\'efinition est muette, 
	              c'est-\`a-dire qu'on peut la remplacer par n'importe quelle autre variable.
	              Ainsi, $\int_{a}^{b} f(x) \dx = \int_{a}^{b} f(t) \dt$. \newline
	              La valeur de $\int_{a}^{b} f(x) \dx$ ne d\'epend pas de la primitive $F$ choisie.
	              En effet, si $F_1$ et $F_2$ sont deux primitives de $f$, alors 
	              il existe $C\in\mathbb{R}$ tel que $F_1 = F_2 + C$.
	              Ainsi, 
	                    \begin{align*}
	                    \left[ F_1 (x) \right]_{a}^{b} & = F_1 (b) - F_1 (a)                    
	                                                     = (F_2 (b) + C ) - (F_2 (a) + C)       \\
	                                                   & = F_2 (b) - F_2 (a)                    
	                                                     = \left[ F_2 (x) \right]_{a}^{b}       \,.
	                    \end{align*}
	              }%
	      \end{defi}
	      %
	      \begin{thm}
	      Soient $f$ et $g$ deux fonction d\'efinies sur un intervalle $I$ et y admettant des primitives.
	      Soient $F$ une primitive de $f$ et $G$ une primitive de $g$.
	      Soient $a$, $b$ et $c$ trois \'el\'ements de $I$ et $\lambda \in \mathbb{R}$.
	            \begin{enumerate}
	                  \item $\int_{a}^{b} f(x) \dx = -\int_{b}^{a} f(x) \dx$
	                  \item $\int_{a}^{b} (f+g) (x) \dx = \int_{a}^{b} f(x) \dx + \int_{a}^{b} g(x) \dx$
	                  \item $\int_{a}^{b} \lambda f(x) \dx = \lambda \int_{a}^{b} f(x) \dx$
	                  \item $\int_{a}^{b} f(x) \dx = \int_{a}^{c} f(x) \dx + \int_{c}^{b} f(x) \dx$
	                  \item Si $f(x) \leq g(x)$ pour tout $x\in [a,b]$ alors 
	                        $\int_{a}^{b} f(x) \dx \leq \int_{a}^{b} g(x) \dx$.            
	            \end{enumerate}
	      \end{thm}
	      %
	      \begin{prop}
	            Soit $f$ une fonction d\'efinie et continue sur un intervalle $[a,b]$ ($a<b$).
	            \begin{enumerate}
	                  \item Si $ f(x) \leq 0$ pour tout $x\in [a,b]$, alors $\int_{a}^{b} f(x) \dx \leq 0$.
	                  \item Si $ f(x) \geq 0$ pour tout $ x\in [a,b]$, alors $\int_{a}^{b} f(x) \dx \geq 0$.
	                  \item $ \left\lvert \int_{a}^{b} f(x) \dx \right\rvert \leq \int_{a}^{b} |f(x)| \dx $.
	                  \item Si $m\leq f(x) \leq M$ pour tout $x\in [a,b]$, 
	                        alors $ m \leq \frac{1}{b-a} \int_{a}^{b} f(x) \dx \leq M$.
	            \end{enumerate}
	      \end{prop}
	      %
	      \begin{thm}[Int\'egration par parties]
	      \index{Integration@Int\'egration!par partie}
	      \index{IPP}
	            Soient $u$ et $v$ deux fonctions d\'erivables sur un intervalle non-vide $I$.
	            Si $u'v$ admet une primitive sur $I$, alors on a 
	                  \begin{equation*}
	                        \int u(x) \cdot v'(x) \dx = u(x) \cdot v(x) - \int u'(x) \cdot v(x) \dx      
	                  \label{eq:IPP fonctions}
	                  \end{equation*}
	            ou, avec diff\'erente notations :
	                  \begin{equation*}
	                        \int u \ \mathrm{d}v = u \cdot v - \int v \ \mathrm{d}u      \,.
	                  \label{eq:IPP differentielles}
	                  \end{equation*}
	            \begin{ex}
	                  On consid\`ere la fonction $f$ d\'efinie par $f(x) = x\cdot \expo{x}$.
	                  On veut trouver la primitive $F$ de $f$ qui v\'erifie $F(1) = 1$.
	                  
	                  Int\'egrons par partie en posant
	                        \begin{align*}
	                              u           & = x       & \text{ et } & &      \mathrm{d}v & = \expo{x} \dx      \\
	                              \mathrm{d}u & = \dx     & \text{    } & &      v           & = \expo{x}          \;.
	                        \end{align*}
	                  On a alors :
	                        \begin{align*}
	                        \int x \cdot \expo{x} \dx & = [x \cdot \expo{x} ] - \int \expo{x} \dx       \\
	                                                  & = x \cdot \expo{x} - \expo{x} + C               \;.
	                        \end{align*}
	                  Posons alors $F(x) = \expo{x} (x-1) + C$.
	                  On cherche alors la valeur de $C$ pour que $F(1) = 1$.
	                  On a $
	                              F(0) = \expo{1} (1-1) + C = C
	                       $,
	                       ce qui nous donne $C = 1$.
	                  D'o\`u $F(x) = \expo{x} (x-1) + 1$.
	            \end{ex}
	      \end{thm}
	      %
	      \begin{thm}
	      \index{Primitives!Changement de variable}
	      Soit $f$ une fonction d\'efinie sur un intervalle $I$ et y admettant une primitive $F$ 
	      et soit $\varphi$ une fonction d\'efinie et d\'erivable sur un intervalle $J$ tel que $I \subseteq \varphi (J)$.
	      La fonction $t \mapsto F\left( \varphi (t) \right)$ est d\'erivable sur $J$ et a pour d\'eriv\'ee 
	      $t \mapsto \varphi'(t) \cdot F' \left( \varphi (t) \right)  = \varphi'(t) \cdot f\left(\varphi (t) \right) $.
	      
	      Si de plus $\varphi$ est une bijection de $J$ vers $I$, on peut poser $x=\varphi (t)$, $x\in I$, on a :
	            \begin{equation*}
	                  \int f(x) \dx = \int f'(\varphi(t)) \cdot \varphi'(t) \dt      \,.
	            \label{eq:changement de variable primitive}
	            \end{equation*}
	            %
	            \begin{ex}
	                  Posons $f(x) = x \sqrt{x+1}$.
	                  On veut trouver une primitive de $f$.
	                  Posons $t = \sqrt{x+1}$, on a $x = t^2 - 1$ et donc $\dx = 2t \dt$.
	                  \begin{align*}
	                        \int x \sqrt{x+1} \dx & = \int (t^2 - 1) \cdot t \cdot 2t \dt                                     
	                                                = \int (2t^4 - 2t^2) \dt                                                  \\
	                                              & = 2 \left[ \frac{t^5}{t} - \frac{t^3}{3} \right] + C                      
	                                                = \frac{2}{5} (x+1)^{\frac{5}{2}} - \frac{2}{3} (x+1)^{\frac{3}{2}} + C   \,.
	                  \end{align*}
	                  Ainsi, la fonction $F$ d\'efinie par 
	                  $F(x) = \frac{2}{5} (x+1)^{\frac{5}{2}} - \frac{2}{3} (x+1)^{\frac{3}{2}}$
	                  est une primitive de $f$.
	            \end{ex}
	      \end{thm}
	      %
      \subsection{Interpr\'etation graphique}
      \index{Integrales@Int\'egrales!Interpretation graphique@Interpr\'etation graphique}
      Soient $f$ une fonction continue et positive sur un intervalle $[a,b]$ ($a<b$)
      et notons par $\mathscr{C}$ sa courbe repr\'esentative dans 
      le plan rapport\'e \`a un rep\`ere orthogonal $(O,\vec{\imath},\vec{\jmath})$.
      L'aire du domaine limit\'e par $\mathscr{C}$, l'axe $x'Ox$ et les droites d'\'equations $x=a$ et $x=b$
      est alors \'egal \`a {$\int_{a}^{b} f(x) \dx$}.
      Cette valeur s'exprime en unit\'e d'aire : c'est-\`a-dire la surface du rectangle de dimensions $1\times1$.

      Notons que si $f$ est n\'egative sur $[a,b]$, alors la valeur de $\int_{a}^{b} f(x) \dx$ est aussi n\'egative.
      
      \noindent
		\begin{tikzpicture}[line cap=round,line join=round,>=stealth,scale=1.2]
			\clip(-2.5,-0.75) rectangle (2.5,2.5);
				\draw [color=black, xstep=.1,ystep=.1, opacity=.2, line width=.2] (-2.75,-0.75) grid (2.75,3.5);
				\draw [color=black, xstep=.5,ystep=.5, opacity=.3, line width=.3] (-2.75,-0.75) grid (2.75,3.5);
				\draw [color=black, xstep=1,ystep=1, opacity=.5, line width=.5] (-2.75,-0.75) grid (2.75,3.5);
				\draw[->,color=black,thick] (-2.75,0) -- (2.5,0);
				\foreach \x in {-2,-1,1,2}
				\draw[shift={(\x,0)},color=black] (0pt,2pt) -- (0pt,-2pt) node[below] {\footnotesize $\x$};
				\draw[->,color=black,thick] (0,-0.75) -- (0,2.5);
				\foreach \y in {,1,2,3}
					\draw[shift={(0,\y)},color=black] (2pt,0pt) -- (-2pt,0pt) node[left] {\footnotesize $\y$};
				\draw[color=black] (0pt,-10pt) node[right] {\footnotesize $0$};
				\draw[smooth,samples=100,domain=-2.75:2.75,thick,loko2] plot(\x,{1/9*((\x)^3-4*(\x)+15)});
				\fill[smooth,samples=100,domain=-1.5:1.5,thick,top color=loko2,bottom color=loko1,opacity=.3] 
					plot(\x,{1/9*((\x)^3-4*(\x)+15)}) -- (1.5,0) -- (-1.5,0) -- cycle;
				\draw[] (0,0.75) node[] {\textcolor{white}{$\int_{a}^{b} f(x) \dx$}};
				\draw[color=loko1!50!black] (-1.5,0) node[] {$\bullet$} node[below] {$a$};
				\draw[color=loko1!50!black] (1.5,0) node[] {$\bullet$} node[below] {$b$};
		\end{tikzpicture}
		%
		\hspace{\stretch{1}}
		%
		\begin{tikzpicture}[line cap=round,line join=round,>=stealth,scale=1.2]
			\clip(-2.5,-2.5) rectangle (2.5,0.75);
				\draw [color=black, xstep=.1,ystep=.1, opacity=.2, line width=.2] (-2.75,-3.5) grid (2.75,0.75);
				\draw [color=black, xstep=.5,ystep=.5, opacity=.3, line width=.3] (-2.75,-3.5) grid (2.75,0.75);
				\draw [color=black, xstep=1,ystep=1, opacity=.5, line width=.5] (-2.75,-3.5) grid (2.75,0.75);
				\draw[->,color=black,thick] (-2.75,0) -- (2.5,0);
				\foreach \x in {-2,-1,1,2}
				\draw[shift={(\x,0)},color=black] (0pt,2pt) -- (0pt,-2pt) node[below] {\footnotesize $\x$};
				\draw[->,color=black,thick] (0,-2.5) -- (0,0.75);
				\foreach \y in {-3,-2,-1}
					\draw[shift={(0,\y)},color=black] (2pt,0pt) -- (-2pt,0pt) node[left] {\footnotesize $\y$};
				\draw[color=black] (0pt,-10pt) node[right] {\footnotesize $0$};
				\draw[smooth,samples=100,domain=-2.75:2.75,thick,loko2] plot(\x,{-1/9*((\x)^3-4*(\x)+15)});
				\fill[smooth,samples=100,domain=-1.5:1.5,thick,top color=loko1,bottom color=loko2,opacity=.3] 
					plot(\x,{-1/9*((\x)^3-4*(\x)+15)}) -- (1.5,0) -- (-1.5,0) -- cycle;
				\draw[] (0,-0.75) node[] {\textcolor{white}{$-\int_{a}^{b} f(x) \dx$}};
				\draw[color=loko1!50!black] (-1.5,0) node[] {$\bullet$} node[above] {$a$};
				\draw[color=loko1!50!black] (1.5,0) node[] {$\bullet$} node[above] {$b$};
		\end{tikzpicture}
		
      Ainsi, dans ce cas, l'aire d\'elimit\'ee par $\mathscr{C}$, l'axe $x'Ox$ et les droites d'\'equations $x=a$ et $x=b$
      est alors \'egal \`a {$-\int_{a}^{b} f(x) \dx$}.
      
      Plus g\'en\'eralement, si la fonction prend \`a la fois des valeurs positives et n\'egatives sur $[a,b]$,
      l'aire du domaine plan correspond \`a $\int_{a}^{b} |f(x)| \dx$.
      La valeur absolue permet de s'assurer qu'on calcule toujours des valeurs positives.
      Pour calculer l'aire du domaine limit\'e par la courbe d'une fonction qui change de signe,
      on peut diviser l'intervalle consid\'er\'e en plusieurs sous-intervalles
      de mani\`ere \`a ce que le signe de la fonction soit constant sur chaque sous-intervalle.
      On consid\`ere alors la valeur absolue des surfaces ainsi calcul\'ees.
      Il suffit ensuite d'additionner les aires trouv\'ees pour obtenir l'aire totale.
      Si on calcule juste la valeur de $\int_{a}^{b} f(x) \dx$, alors 
      la partie n\'egative et la partie positive se soustrairont.
            \begin{figure}[h]
                  \begin{tikzpicture}[>=stealth,x=0.2\actualtextwidth,y=0.2\actualtextwidth]
                        \clip (-2.5,-1.65) rectangle (2.5,1.65);
                        \draw [color=black, xstep=.1,ystep=.1, opacity=.2, line width=.2] (-2.5,-1.65) grid (2.5,1.65);
                        \draw [color=black, xstep=.5,ystep=.5, opacity=.3, line width=.3] (-2.5,-1.65) grid (2.5,1.65);
                        \draw [color=black, xstep=1,ystep=1, opacity=.5, line width=.5] (-2.5,-1.65) grid (2.5,1.65);
                        \draw[->,color=black,thick] (-2.75,0) -- (2.5,0);
                        \foreach \x in {-2,-1,1,2}
                        \draw[shift={(\x,0)},color=black] (0pt,2pt) -- (0pt,-2pt) node[below] {$\x$};
                        \draw[->,color=black,thick] (0,-1.65) -- (0,1.65);
                        \foreach \y in {-1,1}
                        \draw[shift={(0,\y)},color=black] (2pt,0pt) -- (-2pt,0pt) node[left] {$\y$};
                        \draw[color=black] (0pt,-10pt) node[right] {$0$};
                        \fill[samples=250,domain=-1.5:1.5,top color=loko2,bottom color=loko1,opacity=.3] 
                                    plot(\x,{1/3*abs(((\x)^3-4*(\x)))}) -- (1.5,0) -- (-1.5,0) -- cycle;
                        \draw[samples=200,domain=-2.75:2.75,loko2,line width=2.5pt] 
                                    plot(\x,{1/3*abs(((\x)^3-4*(\x)))});
                        \draw[samples=200,domain=-2.75:2.75,white,line width=1.5pt] % Amin'izay hita ilay dashed loko1 line anatiny
                                    plot(\x,{1/3*abs(((\x)^3-4*(\x)))});
                        \draw[smooth,samples=100,domain=-2.75:2.75,loko1, dash pattern=on 5pt off 3pt,line width=.5pt] 
                                    plot(\x,{1/3*((\x)^3-4*(\x))});
                        \draw[color=loko1!50!black] (-1.5,0) node[] {$\bullet$} node[below] {$a$};
                        \draw[color=loko1!50!black] (1.5,0) node[] {$\bullet$} node[below] {$b$};
                        \filldraw[color=loko2,fill=white,line width=2.5pt] 
                             (-2.3,1.25) rectangle node[] {\textcolor{black}{$y=|f(x)|$}} (-1.5,1.55);
                        \draw[color=white,line width=1.5pt] 
                             (-2.3,1.25) rectangle (-1.5,1.55);
                        \filldraw[color=loko1,fill=white, dash pattern=on 5pt off 3pt,line width= .5pt] 
                             (-2.3,-1.25) rectangle node[] {\textcolor{black}{$y=f(x)$}} (-1.5,-1.55);
                  \end{tikzpicture}
            \end{figure}
      
      
      \begin{prop}
            Soient $a>0$ et $f$ une fonction d\'efinie sur $[-a,a]$.
            \begin{itemize}
                  \item Si $f$ est une fonction paire alors $\int_{-a}^{a} f(x) \dx = 2 \cdot \int_{0}^{a} f(x) \dx$.
                  \item Si $f$ est une fonction impaire alors $\int_{-a}^{a} f(x) \dx = 0$.
            \end{itemize}
      \end{prop}
      %
      \begin{thm}
      \index{Integrales@Int\'egrales!Changement de variable}
      Soit $f$ une fonction continue sur $[a,b]$ et $\psi : J \rightarrow [a,b]$ une fonction d\'erivable et bijective.
      
      Si $\psi (\alpha) = a$ et $\psi (\beta) = b$, alors : 
            \begin{equation*}
                  \int_{a}^{b} f(x) \dx = \int_{\alpha}^{\beta} f'(\psi(t)) \cdot \psi'(t) \dt      \,.
            \label{eq:changement de variable integrale}
            \end{equation*}
            %
            \begin{ex}
                  Posons $f(x) = x \sqrt{x+1}$.
                  On veut trouver l'integrale $\int_{0}^{1} f(x) \dx$.
                  Posons $t = \sqrt{x+1}$, on a $x = t^2 - 1$ et donc $\dx = 2t \dt$.
                  Si $x=0$ alors $t=1$ et si $x=1$ alors $t=\sqrt{2}$.
                  \begin{align*}
                  \int_{0}^{1} x \sqrt{x+1} \dx & = \int_{1}^{\sqrt{2}} (t^2 - 1) \cdot t \cdot 2t \dt                      
                                                  = \int_{1}^{\sqrt{2}} (2t^4 - 2t^2) \dt                                   \\
                                                & = 2 \left[ \frac{t^5}{5} - \frac{t^3}{3} \right]_{1}^{\sqrt{2}}           
                                                  = 2 \left[ 
                                                      \left( \frac{4\sqrt{2}}{5}-\frac{2\sqrt{2}}{3} \right) 
                                                      -
                                                      \left( \frac{1}{5} - \frac{1}{3} \right)
                                                      \right]                                                               \\
                                                & = \frac{4}{15} \cdot \left(\sqrt{2} + 1 \right)                           \,.
                  \end{align*}
            \end{ex}
      \end{thm}
      %
      \subsection{Fonction d\'efinie par une int\'egrale}
                  Soit $f$ une fonction d\'efinie et continue sur un intervalle $I$.
            \subsubsection{De la forme $\int_{a}^{x} f(t) \dt$}
                  Comme $f$ est continue, elle admet des primitives sur $I$.
                  Pour tout $a \in I$, on peut d\'efinir la fonction $F$ par \defbox{$F(x) = \int_{a}^{x} f(t) \dt$},
                  qui n'est autre que la primitive de $f$ qui s'annule en $a$.
                  Cette fonction est d\'erivable et a pour d\'eriv\'ee $f$.
            \subsubsection{De la forme $\int_{u(x)}^{v(x)} f(t) \dt$}
                  Soient $u$ et $v$ deux fonctions d\'erivables dans un domaine $\mathcal{D}$, et \`a valeurs dans $I$.
                  La fonction $f$ \'etant continue, elle admet une primitive $F$%
                  \footnote{En fait elle en admet plusieurs, on va juste consid\'erer une d'elles.}
                  
                  On peut alors d\'efinir la fonction $H$ par 
                  \defbox{$H(x) = \int_{u(x)}^{v(x)} f(t) \dt$}.
                  On sait que 
                  \begin{align*}
                        H(x) & = \int_{u(x)}^{v(x)} f(t) \dt = F(v(x)) - F(u(x)) \;.
                  \end{align*}
                  Cette fonction est d\'erivable sur $\mathcal{D}$ et sa d\'eriv\'ee est d\'efinie par
                  \begin{center}
                        \thmbox{
                              \[
                                    H'(x) = v'(x) \cdot f(v(x)) - u'(x) \cdot f(u(x)) .
                              \]
                        }
                  \end{center}









\TCDLessonStyleOff
\clearpage
\providecommand{\Oij}{$(O ; \vec{\imath} , \vec{\jmath})$}
\section{Exercices}
\begin{exo}
      Indiquer les intervalles de $\mathbb{R}$ sur lesquels la fonction $f$ admet des primitives et
      d\'eterminer les primitives de $f$ sur chacun de ces intervalles dans chacun des cas suivants : 
            \begin{enumerate}
                  \item $f : x \mapsto (1+x)^n,\ n \in \mathbb{N}$
                  \item $f : x \mapsto (ax+b)^n,\ n \in \mathbb{N},\ a\in\mathbb{R}^* \text{ et } b\in \mathbb{R}$
                  \item $f : x \mapsto x^2 \cdot \sqrt{1+x^3}$
                  \item $f : x \mapsto \frac{x}{\sqrt{1+x^2}}$
                  \item $f : x \mapsto \left( 1 - \frac{1}{\sqrt[3]{x}} \right)^2$
                  \item $f : x \mapsto \frac{\left(1+\sqrt{x}\right)^2}{\sqrt[3]{x}}$
                  \item $f : x \mapsto 3x^2 (x^2+1) + 2x (x^3 + 1)$
                  \item $f : x \mapsto x^2 \cdot \cos x + 2x \cdot \sin x$
                  \item $f : x \mapsto \frac{x^3 - 3x^2 + 3x -1}{x-1}$
                  \item $f : x \mapsto \frac{x+a}{x^3 + a x^2} ,\ a>0$
                  \item $f : x \mapsto \frac{\sin x - x \cos x}{\sin^2 x}$
                  \item $f : x \mapsto \frac{x \sin x + \cos x}{x^2}$
                  \item $f : x \mapsto 4\sqrt{x^{10} - x^6}$
                  \item $f : x \mapsto \frac{|x-a|}{(x-a)^4}$
                  \item $f : x \mapsto x\cdot \cos \alpha x ,\ \alpha \in \mathbb{R}^*$
                  \item $f : x \mapsto x^2 \cdot \sin \alpha x,\ \alpha \in \mathbb{R}^*$
                  \item $f : x \mapsto \frac{x}{\sqrt{x+1}} \ (\text{poser $t=\sqrt{x+1}$})$
                  \item $f : x \mapsto \sin x \cdot \expo{x}$
                  \item $f : x \mapsto \cos x \cdot \expo{x}$
                  \item $f : x \mapsto \frac{\ln x}{x}$
                  \item $f : x \mapsto \ln x$
                  \item $f : x \mapsto x \cdot \expo{x^2}$
            \end{enumerate}
\end{exo}
%
\begin{exo}
      Calculer les int\'egrales suivantes : 
      \begin{enumerate}
            \item $ \int_{0}^{1} \frac{x}{(1+x^2)^2} \dx $
            \item $ \int_{0}^{2} x^2 \sqrt{1+x^2} \dx $
            \item $ \int_{0}^{4} \frac{x}{\sqrt{1+3x^2}} \dx $
            \item $ \int_{0}^{ \frac{ \sqrt{ \pi } }{2} } \frac{2x}{ \cos ^2 x^2 } \dx $
            \item $ \int_{-2}^{3} |x| \dx $
            \item $ \int_{-1}^{2} x \cdot |x| \dx $
            \item $ \int_{-1}^{2} |x| \cdot \expo{2x^2} \dx $
            \item $ \int_{0}^{1} x^2 \cdot \expo{x} \dx $
            \item $ \int_{1}^{e} x \cdot \ln x \dx  $
            \item $ \int_{-2}^{2} \frac{x^2 + 2x + 1}{x+1} \dx $
            \item $ \int_{-2}^{2} \frac{x^2 + 2x + 1}{x^2 + 1} \dx $
            \item $ \int_{0}^{1} \frac{2x^2 + 3x -2}{x^2-1} \dx $
            \item $ \int_{0}^{1} \frac{x^2 - 3x + 15}{x^2 -5x +6} \dx $ 
                  (Trouver $a$, $b$ et $c$ tels que
                  $ \frac{x^2 - 3x + 15}{x^2 -5x +6} = a+\frac{b}{x-2}+\frac{c}{x-3} $
                  )
      \end{enumerate}
\end{exo}
%
\begin{exo}
      Calculer les limites suivantes : 
      \begin{enumerate}
            \item $\displaystyle{ \lim_{x\rightarrow 0} \int_{x}^{1} \frac{1}{\sqrt{t}} \dt}$
            \item $\displaystyle{ \lim_{x\rightarrow 0} \int_{x}^{1} \frac{1}{t^{\alpha}} \dt}$ o\`u $0 < \alpha < 1$.
            \item $\displaystyle{ \lim_{x\rightarrow 1} \int_{0}^{x} \frac{t}{\sqrt{1-t^2}} \dt}$
            \item $\displaystyle{ \lim_{x\rightarrow +\infty} \int_{1}^{x} \frac{1}{x} \dt}$
            \item $\displaystyle{ \lim_{x\rightarrow +\infty} \int_{1}^{x} \frac{1}{x^{\alpha}} \dt}$ o\`u $\alpha > 1$.
            \item $\displaystyle{ \lim_{x\rightarrow +\infty} \int_{0}^{x} \expo{-x} \dt}$
      \end{enumerate}
\end{exo}
%
\begin{exo}
      On consid\`ere les fonction $f$ et $g$ d\'efinies sur $[1,+\infty[$ respectivement par :
            \[
                  f(x) = \frac{1}{x} \text{ et } g(x) = \frac{1}{x^2}      .
            \]
      On note $\mathscr{C}$ et $\Gamma$ leurs courbes repr\'esentatives respectives.
      On d\'esigne par $\mathscr{A}(\lambda)$ l'aire du plan compris entre 
      la courbe $\mathscr{C}$, $\Gamma$ et la droite d'\'equation $x=\lambda$ ($\lambda > 1$).
      \begin{figure}[H]
            \begin{tikzpicture}[line cap=round,line join=round,>=stealth,x=1.0cm,y=1.0cm]
            \clip(-0.5,-0.5) rectangle (11,2.5);
            \draw[->,color=black] (-0.25,0) -- (11,0);
            \foreach \x in {1,2,3,4,5,6,7,8,9,10,11}
            \draw[shift={(\x,0)},color=black] (0pt,2pt) -- (0pt,-2pt) node[below] {\footnotesize $\x$};
            \draw[->,color=black] (0,-0.5) -- (0,2.5);
            \foreach \y in {1,2}
            \draw[shift={(0,\y)},color=black] (2pt,0pt) -- (-2pt,0pt) node[left] {\footnotesize $\y$};
            \draw[color=black] (0pt,-10pt) node[right] {\footnotesize $0$};
            \fill[left color=loko2,right color=loko1,samples=100,domain=1:7,opacity=.5] 
                  plot(8-\x,{1/(8-\x)}) -- plot(\x,{1/(\x)^2}) -- cycle;
            \draw[color=loko2, smooth,samples=100,domain=1:12.5] plot(\x,{1/(\x)});
            \draw[color=loko1, smooth,samples=100,domain=1:12.5] plot(\x,{1/(\x)^2});
            \draw[color=loko2, smooth,samples=100,domain=.1:1, dash pattern=on 2pt off 2pt] plot(\x,{1/(\x)});
            \draw[color=loko1, smooth,samples=100,domain=.1:1, dash pattern=on 2pt off 2pt] plot(\x,{1/(\x)^2});
            \end{tikzpicture}
      \end{figure}
      Exprimer $\mathscr{A}(\lambda)$ en fonction de $\lambda$.
      Calculer $\lim_{\lambda \rightarrow +\infty} \mathscr{A} (\lambda)$.
\end{exo}
%
\begin{exo}
      On consid\'ere la suite $(S_n)_{n\in\mathbb{N}^*}$ d\'efinie par :
            $\displaystyle{ S_n = \sum_{i=1}^{n} \frac{1}{i^2} }$.
       Pour tout $n\in\mathbb{N}^*$, on pose $ \displaystyle{I_n = \int_{1}^{n} \frac{\dx}{x^2}  } $.
      \begin{enumerate}
            \item Montrer que les suites $(S_n)$ et $(I_n)$ sont croissantes.
            \item \begin{enumerate}
                        \item Interpreter $I_n$ graphiquement.
                        \item Que peut \^etre la signification graphique de $S_n$\,?
                  \end{enumerate}
            \item \begin{enumerate}
                        \item Calculer $\lim_{n\rightarrow+\infty} I_n$. 
                              En d\'eduire un majorant de $(I_n)$.
                        \item Justifier que $S_n \leq S_1 + I_n$ pour tout $n \geq 2$.
                        \item Que peut-on alors dire de la suite $(S_n)$\,?
                  \end{enumerate}
      \end{enumerate}
      %
	            \begin{tikzpicture}[line cap=round,line join=round,>=stealth,x=0.145\linewidth,y=0.145\linewidth]
            \clip(-0.2,-0.2) rectangle (6.25,2.1);
            \draw [color=black, xstep=1,ystep=1, line width=0.5, opacity=0.5] (-0.18,-0.21) grid (6.23,2.08);
            \draw [color=black, xstep=.5,ystep=.5, line width=0.3, opacity=0.3] (-0.18,-0.21) grid (6.23,2.08);
            \draw [color=black, xstep=.1,ystep=.1, line width=0.2, opacity=0.2] (-0.18,-0.21) grid (6.23,2.08);
            \draw[smooth,samples=100,domain=.1:6.25,pattern color=loko1!50!black,pattern=north west lines,opacity=.75] 
                     plot(\x,{1/(\x)^2}) -- (6.25,0) -- (0,0) -- cycle;
            \begin{scope}
                  \foreach \x in {1,2,3,4,5,6,7}
                  \filldraw[color=loko2!50!black,fill=loko2,opacity=.75] 
                           (\x-1,0) rectangle 
                           %node[opacity=1,above] {\large{$\boldmath{\frac{1}{\x ^2}}$}} 
                           (\x,{1/(\x)^2});
            \end{scope}
            \draw[->,color=black,thick] (-0.18,0) -- (6.23,0);
            \draw[->,color=black,thick] (0,-0.21) -- (0,2.08);
            \end{tikzpicture}
\end{exo}
%
\begin{exo}
      Soit $(I_n)_{n\in\mathbb{N}^*}$ la suite d\'efinie par 
            $\displaystyle{
                  I_n = \int_{0}^{n} \frac{1}{x^3 + 1} \dx 
            }$.
      \begin{enumerate}
            \item Montrer que la suite $(I_n)$ est croissante.
            \item Montrer que pour tout $n\in\mathbb{N}^*$, on a 
                  $\displaystyle{
                        I_n \leq I_1 + \int_{1}^{n} \frac{\dx}{x^3}
                  }$.
            \item En d\'eduire que $(I_n)$ est major\'ee. Que peut-on en conclure\,?
      \end{enumerate}
\end{exo}
%
\begin{exo}
      On d\'efinit la suite $(I_n)_{n\in\mathbb{N}*}$ par 
            \[
                  I_n = \int_{0}^{n} \frac{\sqrt{x}}{x^4+1} \dx       .
            \]
      \begin{enumerate}
            \item Montrer que $(I_n)$ est croissante.
            \item Montrer que pour tout $n\geq 2$, on a 
                  \[
                        I_n \leq I_1 + \int_{1}^{n} \frac{1}{x^2} \dx      .
                  \]
            \item D\'eduire des questions pr\'ec\'edentes que $(I_n)$ est convergente.
      \end{enumerate}
\end{exo}
%
\begin{exo}
      Soit $\alpha$ un nombre entier naturel sup\'erieur ou \'egal \`a $2$.
      On d\'efinit la suite $(S_n)_{n\in\mathbb{N}}$ par 
            \[
                  S_n= \sum_{i=1}^n \frac{1}{i^2}      .
            \]
      \begin{enumerate}
            \item On d\'efinit la suite $(I_n)_{n\in\mathbb{N}*}$ par 
                  \[
                        I_n = \int_{1}^{n} \frac{1}{x^\alpha} \dx      .
                  \]
                  \begin{enumerate}
                        \item Montrer que $(I_n)_{n\in\mathbb{N}*}$ est croissante.
                        \item Montrer que pour tout $n\geq 2$.
                              \[
                                    S_n \leq S_1+ I_n      .
                              \]
                  \end{enumerate}
            \item \begin{enumerate}
                        \item Montrer que $(S_n)_{n\in\mathbb{N}*}$ est croissante.
                        \item En majorant, d'abord $(I_n)_{n\in\mathbb{N}*}$, 
                              montrer que $(S_n)_{n\in\mathbb{N}*}$ est aussi major\'ee.
                              Conclure.
                  \end{enumerate}
      \end{enumerate}
\end{exo}
%%%%%%%%%%%%%%%%%%%%%%%%%%%%%%%%%%%%%%%%%%%%%%%%%%%%%%%%%%%%%%%%%%%%%%%%%%%%%%%%%%%
\begin{exo} ~ \newline
\begin{enumerate}
\item Calculer la fonction d\'eriv\'ee de la fonction $x \mapsto \ln \left( x + \sqrt{x^2 + 2} \right) $.
      En d\'eduire une primitive de la fonction $f$ d\'efinie par $f(x) = \frac{1}{\sqrt{x^2 + 2}}$.
\item On consid\`ere les fonctions $g$ et $h$ d\'efinies sur $\mathbb{R}$ par 
      $g(x) = \frac{x^2}{\sqrt{x^2 + 2}}$ et $h(x) = \sqrt{x^2 + 2}$.
      \begin{enumerate}
      \item V\'erifier que $h - g = 2 \; f$ et que $h + g$ est la d\'eriv\'ee de $x \mapsto x \; \sqrt{x^2 + 2}$.
      \item En d\'eduire une primitive de $g$ et $h$ sur $\mathbb{R}$.
      \end{enumerate}
\end{enumerate}
\end{exo}
%
%%%%%%%%%%%%%%%%%%%%%%%%%%%%%%%%%%%%%%%%%%%%%%%%%%%%%%%%%%%%%%%%%%%%%%%%%%%%%%%%%%%
\begin{exo}
            Soit $(U_n)$ la suite d\'efinie, pour tout entier naturel $n>0$, par
            \[
                  U_n=\int_{0}^{1}x^n\ln(x+1)\dx.
            \]
            \begin{enumerate}
                  \item Etudier le signe de $(x^{n+1} - x^n)$ sur $[0 ; 1]$.
                  \item En d\'eduire le sens de variation de $(U_n)$.
                  \item La suite $(U_n)$ est-elle convergente\,? Justifier.
                  \item D\'emontrer que pour tout entier naturel $n>0$, on a $0 \leq U_n \leq \frac{\ln 2}{n+1}$.
                  \item En d\'eduire $\displaystyle{\lim_{n\rightarrow+\infty} U_n}$.
            \end{enumerate}
\end{exo}
%%%%%%%%%%%%%%%%%%%%%%%%%%%%%%%%%%%%%%%%%%%%%%%%%%%%%%%%%%%%%%%%%%%%%%%%%%%%%%%%%%%




















